\documentclass[11pt,a4paper]{article}
\usepackage[margin=1in]{geometry}
\usepackage[T1]{fontenc}
\usepackage[utf8]{inputenc}
\usepackage{lmodern}
\usepackage{hyperref}
\usepackage{enumitem}
\usepackage{xcolor}
\usepackage{listings}

\hypersetup{
  colorlinks=true,
  linkcolor=blue,
  urlcolor=blue
}

\lstset{
  basicstyle=\ttfamily\small,
  breaklines=true,
  frame=single,
  columns=fullflexible
}

\title{Lab Logbook: Open Authorization Protocol Project}
\author{Group: \texttt{<fill group members>}\\Course: 02244 Logic for Security}
\date{Hand-in Date: \texttt{<fill date>}}

\begin{document}
\maketitle

\section*{How this logbook is organized}
This logbook is written as a continuous technical narrative, week by week, in the same order as the course development plan.
For every week, the text explicitly addresses the four points required in the project handout:
what was changed and why, which modeling assumptions and goals were used, what problems or attacks were encountered, and what design decisions were taken after analysis.
The purpose of this format is to provide a report-ready story of protocol evolution rather than a list of isolated notes.

Please replace each \texttt{Author} field below with one group member name to satisfy the individualized authorship requirement.

\section{Week 2: Dolev--Yao Baseline}
\textbf{Author: \texttt{<name>}}

In Week 2 we deliberately built a minimal baseline protocol in order to establish a working formal model before introducing realistic trust assumptions.
The protocol captures the core open-authorization intent: owner \texttt{A} delegates photo access at server \texttt{B} to service \texttt{P}.
At this stage, all public keys are assumed to be globally known and there is no identity-provider interaction yet.
The selected scope and validity are represented by symbolic fields (\texttt{F,T}), and the initial goals are secrecy of photo data and authorization consistency at the resource server.

The key technical result of Week 2 is that this baseline is insecure under Dolev--Yao reasoning.
Using the passive-observer analysis required by the assignment, we showed that the protocol structure allows derivation of the protected photo data when the delegate role is intruder-controlled.
This outcome is not a failure of the process; it is the expected starting point that justifies the subsequent protocol refinements in Weeks 3--6.
Related artifacts for this week are:
\texttt{Week2\_Problem1\_Informal\_Protocol\_Outline.md},
\texttt{Week2\_Problem2\_First\_AnB\_Version.anb},
and \texttt{Week2\_Problem3\_Static\_DolevYao\_Analysis.md}.

\section{Week 3: Identity Provider + Lazy Intruder Task}
\textbf{Author: \texttt{<name>}}

Week 3 introduces the assignment's central trust mechanism: delegation via a trusted identity provider \texttt{I}.
The main protocol (\texttt{Week3\_Problem1\_AnB\_With\_IdentityProvider.anb}) extends the baseline by requiring \texttt{A} to prove knowledge of \texttt{pw(A,I)} through a hash-based challenge response, after which \texttt{I} issues a signed grant token.
This design explicitly follows the requirement that the password is not used as an encryption key.

From a modeling perspective, this week transitions from ``direct owner signature'' assumptions to ``IDP-issued authorization'' assumptions.
The symbolic goals remain secrecy and authentication oriented, but the trust source is now explicit and centralized in \texttt{I}.
OFMC runs show that security issues persist, including secrecy or weak-auth failures depending on role instantiation and refinement choices.
An additional refinement with fixed honest provider (\texttt{Week3\_Problem3\_Refined\_HonestProvider.anb}) was kept to document this behavior.

The Week 3 special task was completed in full using the lecture method.
The file \texttt{Week3\_Problem2\_LazyIntruder\_Executability\_of\_P.md} gives a formal, ordered lazy-intruder derivation showing that role \texttt{P} is executable under \texttt{P=i} with honest \texttt{A,B,I}.

\section{Week 4: Formats, Key Lookup, Type-Flaw Resistance}
\textbf{Author: \texttt{<name>}}

Week 4 focuses on typing discipline and implementation realism.
The main protocol (\texttt{Week4\_Problem1\_MainProtocol\_Formats.anb}) replaces ambiguous concatenation-heavy structures with explicit format constructors (\texttt{fReq}, \texttt{fAuthProof}, \texttt{fGrant}, \texttt{fUse}, \texttt{fData}).
In parallel, we introduced a separate key-lookup protocol (\texttt{Week4\_Problem2\_AnB\_KeyLookup\_Protocol.anb}) so that \texttt{A} obtains public-key statements from \texttt{I} instead of assuming universal key knowledge by default.

This week also required practical debugging work.
A fully wrapped format style initially caused multiple non-executability errors in OFMC because some terms became hard to derive under role knowledge constraints.
We therefore simplified selected message boundaries while preserving format separation where it matters for type-flaw analysis.
The resulting models are executable and continue to expose weak-auth issues, which are retained as evidence in the development narrative.

The formal Week 4 special-task argument is documented in
\texttt{Week4\_Problem3\_TypeFlawResistance\_Check.md}.
Its core claim is that distinct top-level format constructors prevent cross-type unification except where message kinds are intentionally identical.

\section{Week 5: Pseudonymous Channels (TLS-style abstraction)}
\textbf{Author: \texttt{<name>}}

In Week 5 we produced a channel-centric protocol variant
(\texttt{Week5\_Problem1\_PseudonymousChannels.anb}) using pseudonymous channel notation, as required by the lecture's TLS abstraction discussion.
The design intention was to move from explicit cryptographic transport details toward communication assumptions that better resemble real deployments with server-authenticated TLS semantics.
Business logic was kept as close as possible to Week 4 so changes could be attributed to channel abstraction rather than message intent.

The model is executable, but OFMC returns \texttt{ATTACK\_FOUND} for secrecy with leaked \texttt{fData(Photos(1))}, documented in
\texttt{Week5\_Problem1\_PseudonymousChannels\_Result.md}.
This week therefore demonstrates that replacing low-level crypto syntax with channels improves clarity and modularity, but does not by itself guarantee secure authorization behavior.

\section{Week 6: Guessable Password Security Task}
\textbf{Author: \texttt{<name>}}

Week 6 addresses the special task on low-entropy password robustness.
We first modeled \texttt{pw(A,I)} as a guessable secret in
\texttt{Week6\_Problem1\_GuessablePassword\_Check.anb}; OFMC reported \texttt{ATTACK\_FOUND} with goal \texttt{guesswhat}, confirming an offline-guessing vulnerability.
We then introduced a hardened proof structure in
\texttt{Week6\_Problem2\_GuessablePassword\_Hardened.anb} by mixing an additional fresh value \texttt{K} into the password-hash witness.
This removed \texttt{guesswhat} as the first failing property, but secrecy attacks remained.

To isolate the assignment question itself, we created a dedicated model
(\texttt{Week6\_Problem3\_GuessablePassword\_OnlyCheck.anb}) with only the guessable-secret goal active.
With bounded sessions (\texttt{--numSess 1}), OFMC returned \texttt{NO\_ATTACK\_FOUND}.
Therefore, within the tested bound and for the dedicated guessing property, the Week 6 objective is satisfied.
All raw outcomes are summarized in \texttt{Week6\_Problem1\_Result\_Summary.md}.

\section{Reproducibility commands (PowerShell)}
\textbf{Author: \texttt{<name>}}

\begin{lstlisting}
# Week 2
ofmc --numSess 2 "ProjectProtocol\Week2_Problem2_First_AnB_Version.anb"

# Week 3
ofmc --numSess 2 "ProjectProtocol\Week3_Problem1_AnB_With_IdentityProvider.anb"
ofmc --numSess 2 "ProjectProtocol\Week3_Problem3_Refined_HonestProvider.anb"

# Week 4
ofmc --numSess 2 "ProjectProtocol\Week4_Problem1_MainProtocol_Formats.anb"
ofmc --numSess 2 "ProjectProtocol\Week4_Problem2_AnB_KeyLookup_Protocol.anb"

# Week 5
ofmc --numSess 2 "ProjectProtocol\Week5_Problem1_PseudonymousChannels.anb"

# Week 6
ofmc --numSess 2 "ProjectProtocol\Week6_Problem1_GuessablePassword_Check.anb"
ofmc --numSess 2 "ProjectProtocol\Week6_Problem2_GuessablePassword_Hardened.anb"
ofmc --numSess 1 "ProjectProtocol\Week6_Problem3_GuessablePassword_OnlyCheck.anb"
\end{lstlisting}

\section{Resources and collaboration disclosure}
\textbf{Author: \texttt{<name>}}

\subsection*{Tools and materials used}
\begin{itemize}[leftmargin=1.2em]
  \item Course slides (Weeks 2--6 topics: Dolev--Yao, lazy intruder, formats, channels, privacy/guessing context).
  \item OFMC 2024 command-line model checker.
  \item Assignment handout \texttt{protocol-project (2).pdf}.
\end{itemize}

\subsection*{AI usage disclosure}
AI assistant support was used for:
\begin{itemize}[leftmargin=1.2em]
  \item editing and restructuring AnB snapshots,
  \item formatting notes into weekly logbook entries,
  \item drafting presentation/logbook text.
\end{itemize}
All protocol designs, assumptions, and final submitted claims were manually reviewed and validated against OFMC outputs.

\end{document}
